\chapter{Tabular Data Analysis}
\label{chap:tabular}
% Tabular chapter, owner: Trần Lê Hoàng Long

\section{Problem and Objectives}
\textbf{Problem:} [Classification or regression task and its real-world goal.]

\section{Dataset}
\begin{itemize}[leftmargin=1.5cm]
    \item Name: [Fill in]
    \item Source and URL: [Fill in]
    \item Rows / columns: [Fill in]
    \item Target: [Fill in]
    \item Categorical / numerical columns: [Fill in]
    \item Missing values: [Count, ratio and affected columns]
    \item License / citation: [Fill in]
\end{itemize}

\section{Exploratory Data Analysis}
% Numeric distributions, categorical frequencies, correlation heatmap, missing-value chart
[Key findings; save figures to figures/.]

\section{Data Preprocessing}
\subsection{Data Quality Checks}
[Duplicates, wrong dtypes, invalid values, missing values.]

\subsection{Comparing Missing-Value Strategies}
[Drop rows, mean/median, mode or KNNImputer; selection criterion.]

\subsection{Outlier Handling}
[IQR or z-score; before vs after.]

\subsection{Encoding and Scaling}
[One-Hot/Ordinal encoding; StandardScaler/MinMaxScaler/RobustScaler.]

\subsection{Pipeline and Data Leakage Prevention}
[Train/validation/test split, Pipeline and ColumnTransformer.]

\section{Models and Experimental Setup}
[Baseline, Logistic Regression/Decision Tree/Random Forest or other suitable models,
data split and tuning on validation.]

\section{Results and Evaluation}
% Imbalanced data: also report F1, recall, ROC-AUC
\begin{table}[H]
\centering
\caption{Comparison of tabular models}
\begin{tabular}{lrrrr}
\toprule
Model & Accuracy & Precision & Recall & F1-score\\
\midrule
Baseline & -- & -- & -- & --\\
Model 2 & -- & -- & -- & --\\
Best model & -- & -- & -- & --\\
\bottomrule
\end{tabular}
\end{table}

\section{Question-Driven Figure}
\textbf{Question:} [E.g. Does handling class imbalance improve minority-class recall?]

\begin{figure}[H]
\centering
\figureplaceholder{tabular_creative.png}{Chart with numbers to compare.}
\caption{[State the conclusion with numbers matching the notebook.]}
\label{fig:tabular-creative}
\end{figure}

\section{Limitations and Future Work}
[Dataset limits, bias risks, model limitations, possible extensions.]

\section{References}
[Dataset source, model documentation and other sources used.]
